\paragraph{Problem 1 answer}\GradescopeComment{This should be on page 3 (top) of your PDF.}

\begin{enumerate}[label=(\alph*)]

\item[(a)]  Here's a comma-separated list of the nodes in the topological order as computed by the algorithm:

  \begin{blue}
    \[
      \REPLACEME
    \]
  \end{blue}

\item[(b)]

  \begin{blue}
  \(
    \begin{array}{r@{:~~}l}
      \text{node} & \text{parent} \\ \hline

      a & \REPLACEME \\

      b & \REPLACEME \\

      c & \REPLACEME \\

      d & \REPLACEME \\

      e & \REPLACEME \\

      f & \REPLACEME \\

      g & \REPLACEME \\

      h & \REPLACEME \\

      i & \REPLACEME \\

      j & \REPLACEME
    \end{array}
  \)
\end{blue}
  
\item[(c)]  Here are comma-separated lists of the forward, back, and cross edges. \smallskip

  \begin{tabular}{rl}

    Forward edges:
    & \BLUE{\(
      \REPLACEME
      \)}

    \\

    Back  edges:
    & \BLUE{\(
      \REPLACEME
      \)}

    \\

    Cross edges:
    & \BLUE{\(
      \REPLACEME 
      \)}
  \end{tabular}
  
\end{enumerate}

\newpage

\paragraph{Problem 1(d) answer (OPTIONAL!)}\GradescopeComment{This should be on page 4 (top) of your PDF.}

\begin{blue}
  
% HERE ARE TWO WAYS TO DRAW A ROOTED TREE IN LaTeX.
% UNCOMMENT ONE OF THEM AND EDIT IT TO MAKE YOUR ANSWER.

%%%%%%%%%%%%%%%%%%%%%%%%%%%%%%%%%%%%%%%%%%%%%%%%%%%%%%
%%% METHOD 1

% \begin{center}
%   \begin{tikzpicture}[
%     every node/.style={rounded rectangle, draw},
%     every path/.style={->}
%     ]
%     \node (a) [label={180:11}, label={0:16}] {a}; 
%     \node (b) [below =of a, label={180:14}, label={0:15}] {b}; 
%     \node (c) [below =of b, label={180:3}, label={0:6}] {c}; 

%     \node (d) [left =of a, xshift=-20pt, label={180:1}, label={0:10}] {d}; 
%     \node (e) [below =of d, label={180:2}, label={0:9}] {e}; 
%     \node (f) [below =of e, xshift=-10pt, label={180:7}, label={0:8}] {f}; 
%     \node (g) [below =of f, xshift=10pt, label={180:4}, label={0:5}] {g}; 

%     \node (h) [right =of b, xshift=20pt, label={180:12}, label={0:13}] {h}; 
%     \node (i) [above =of h, label={180:19}, label={0:20}] {i}; 

%     \node (j) [right =of i, xshift=20pt, label={180:17}, label={0:18}] {j}; 

%     \draw (a) -- (b); 
%     \draw (b) -- (c); 
%     \draw (d) -- (e); 
%     \draw (e) -- (f); 
%     \draw (f) -- (g); 
%     \draw (e) -- (c); 
%     \draw (c) -- (g); 
%     \draw (a) -- (h); 
%     \draw (i) -- (h); 
%     \draw (b) -- (h); 
%   \end{tikzpicture}
% \end{center}

%%%%%%%%%%%%%%%%%%%%%%%%%%%%%%%%%%%%%%%%%%%%%%%%%%%%%%
%%%   METHOD 2
  
%  % DRAW IT ON PAPER, SCAN IT TO A PDF, THEN INCLUDE THE PDF IN YOUR DOCUMENT AS FOLLOWS:

%   \includegraphics[height=3in,width=6in]{hwk_demo_scan}
  
%   % SEE https://texblog.org/2012/02/23/crop-figures-with-includegraphics/ IF YOU WANT TO CROP THE INCLUDED IMAGE.
  
\end{blue}

\vfill \vfill

\paragraph{Problem \arabic{problem} collaboration and citations}~\GradescopeComment{This should be on page 4 (bottom) of your PDF.}

I collaborated on this problem with the following people:
\begin{blue}
  \begin{itemize}
  \item \REPLACEME                    % write "none" if none
  \end{itemize}
\end{blue}

My answers use ideas from the following sources (other than the lecture notes and textbooks): 
\begin{blue}
  \begin{itemize}
  \item \REPLACEME                    % write "none" if none
  \end{itemize}
\end{blue}



%%% Local Variables:
%%% mode: latex
%%% TeX-master: "homework_6"
%%% End:
